% Options for packages loaded elsewhere
\PassOptionsToPackage{unicode}{hyperref}
\PassOptionsToPackage{hyphens}{url}
\PassOptionsToPackage{dvipsnames,svgnames,x11names}{xcolor}
%
\documentclass[
  11pt,
  a4paper,
]{article}
\usepackage{amsmath,amssymb}
\usepackage{iftex}
\ifPDFTeX
  \usepackage[T1]{fontenc}
  \usepackage[utf8]{inputenc}
  \usepackage{textcomp} % provide euro and other symbols
\else % if luatex or xetex
  \usepackage{unicode-math} % this also loads fontspec
  \defaultfontfeatures{Scale=MatchLowercase}
  \defaultfontfeatures[\rmfamily]{Ligatures=TeX,Scale=1}
\fi
\usepackage{lmodern}
\ifPDFTeX\else
  % xetex/luatex font selection
  \setmainfont[]{Latin Modern Roman}
  \setsansfont[]{Latin Modern Sans}
  \setmonofont[Scale=0.85]{DejaVu Sans Mono}
\fi
% Use upquote if available, for straight quotes in verbatim environments
\IfFileExists{upquote.sty}{\usepackage{upquote}}{}
\IfFileExists{microtype.sty}{% use microtype if available
  \usepackage[]{microtype}
  \UseMicrotypeSet[protrusion]{basicmath} % disable protrusion for tt fonts
}{}
\usepackage{xcolor}
\usepackage[top=23mm,bottom=23mm,left=23mm,right=23mm]{geometry}
\usepackage{longtable,booktabs,array}
\usepackage{calc} % for calculating minipage widths
% Correct order of tables after \paragraph or \subparagraph
\usepackage{etoolbox}
\makeatletter
\patchcmd\longtable{\par}{\if@noskipsec\mbox{}\fi\par}{}{}
\makeatother
% Allow footnotes in longtable head/foot
\IfFileExists{footnotehyper.sty}{\usepackage{footnotehyper}}{\usepackage{footnote}}
\makesavenoteenv{longtable}
\setlength{\emergencystretch}{3em} % prevent overfull lines
\providecommand{\tightlist}{%
  \setlength{\itemsep}{0pt}\setlength{\parskip}{0pt}}
\setcounter{secnumdepth}{-\maxdimen} % remove section numbering
% definitions for citeproc citations
\NewDocumentCommand\citeproctext{}{}
\NewDocumentCommand\citeproc{mm}{%
  \begingroup\def\citeproctext{#2}\cite{#1}\endgroup}
\makeatletter
 % allow citations to break across lines
 \let\@cite@ofmt\@firstofone
 % avoid brackets around text for \cite:
 \def\@biblabel#1{}
 \def\@cite#1#2{{#1\if@tempswa , #2\fi}}
\makeatother
\newlength{\cslhangindent}
\setlength{\cslhangindent}{1.5em}
\newlength{\csllabelwidth}
\setlength{\csllabelwidth}{3em}
\newenvironment{CSLReferences}[2] % #1 hanging-indent, #2 entry-spacing
 {\begin{list}{}{%
  \setlength{\itemindent}{0pt}
  \setlength{\leftmargin}{0pt}
  \setlength{\parsep}{0pt}
  % turn on hanging indent if param 1 is 1
  \ifodd #1
   \setlength{\leftmargin}{\cslhangindent}
   \setlength{\itemindent}{-1\cslhangindent}
  \fi
  % set entry spacing
  \setlength{\itemsep}{#2\baselineskip}}}
 {\end{list}}
\usepackage{calc}
\newcommand{\CSLBlock}[1]{\hfill\break\parbox[t]{\linewidth}{\strut\ignorespaces#1\strut}}
\newcommand{\CSLLeftMargin}[1]{\parbox[t]{\csllabelwidth}{\strut#1\strut}}
\newcommand{\CSLRightInline}[1]{\parbox[t]{\linewidth - \csllabelwidth}{\strut#1\strut}}
\newcommand{\CSLIndent}[1]{\hspace{\cslhangindent}#1}
\ifLuaTeX
\usepackage[bidi=basic]{babel}
\else
\usepackage[bidi=default]{babel}
\fi
\babelprovide[main,import]{british}
\ifPDFTeX
\else
\babelfont{rm}[]{Latin Modern Roman}
\fi
% get rid of language-specific shorthands (see #6817):
\let\LanguageShortHands\languageshorthands
\def\languageshorthands#1{}
\usepackage{fancyhdr}
\usepackage{fvextra}
\usepackage{etoolbox}
\usepackage{microtype}
\usepackage{needspace}
\pagestyle{fancy}
\fancyhf{}
\fancyhead[L]{\small Evidence verification after compaction}
\fancyhead[R]{\small Working draft}
\fancyfoot[C]{\thepage}
\setlength{\headheight}{14pt}
\setlength{\emergencystretch}{2em}
\AtBeginEnvironment{longtable}{\small\setlength{\tabcolsep}{4pt}}
\AtBeginEnvironment{CSLReferences}{\small\raggedright}
\RecustomVerbatimEnvironment{verbatim}{Verbatim}{fontsize=\small,breaklines=true,breakanywhere=true}
\ifLuaTeX
  \usepackage{selnolig}  % disable illegal ligatures
\fi
\usepackage{bookmark}
\IfFileExists{xurl.sty}{\usepackage{xurl}}{} % add URL line breaks if available
\urlstyle{same}
\hypersetup{
  pdftitle={When to Reopen the Source: Evidence Verification after Native Agent Compaction},
  pdflang={en-GB},
  colorlinks=true,
  linkcolor={blue},
  filecolor={Maroon},
  citecolor={black},
  urlcolor={blue},
  pdfcreator={LaTeX via pandoc}}

\title{When to Reopen the Source: Evidence Verification after Native
Agent Compaction}
\usepackage{etoolbox}
\makeatletter
\providecommand{\subtitle}[1]{% add subtitle to \maketitle
  \apptocmd{\@title}{\par {\large #1 \par}}{}{}
}
\makeatother
\subtitle{Methods and protocol working draft}
\author{}
\date{28 November 2026}

\begin{document}
\maketitle

\textbf{Status.} This expanded research draft specifies a proposed
method and an evaluation protocol. It is not formatted as an eight-page
conference submission. It is not a completed empirical paper, a
preregistration, or a claim that the complete proposed system has
already been implemented. Its model-based performance observations come
from an earlier synthetic development study; the freshness
counterexample, engineering checks, host-integration attempts, and
source/environment qualification are discussed separately. No new paired
method comparison, confirmatory estimate, or completed deployment is
asserted. Numerical result cells for the proposed study are deliberately
unfilled. Human authors, affiliations, contributions, and the final
submission format remain to be confirmed.

\section{Abstract}\label{abstract}

An agent continuing a task after native context compaction may retain a
conclusion while losing the evidence needed to determine whether that
conclusion still applies. Reopening raw history can recover missing
information, but a relevant passage may describe an earlier code
revision, an incomplete test run, or an authorization that was later
withdrawn. This working draft studies whether explicit verification
improves subsequent decisions beyond strong raw-history retrieval. We
specify a selective verification layer that discovers candidate claims
from observable task state, retrieves and opens original sources without
evaluator-provided locations, and checks type-specific requirements for
test status, current versions, and historical authorization. Reuse of
previously verified evidence requires a certificate bound to a complete
archive view, a current archive watermark, the relevant object state,
and a valid time interval; a recent timestamp alone is insufficient. The
proposed evaluation pairs native continuations under a shared model,
archive, and evidence budget, comparing the method with Raw-RAG,
always-on verification, and reproducible neighboring approaches. Its
primary endpoint is task-level evidence-supported decision success. Cost
claims require both quality noninferiority and lower total standardized
processing cost. The draft makes no claim of empirical superiority: it
provides the method specification, falsifiable hypotheses, and evidence
required to determine when source verification is useful.

\clearpage

\section{1. Introduction}\label{introduction}

A coding agent runs tests, changes a file, receives an updated
requirement, and continues after its host compacts the conversation. The
remaining context says that the tests passed. That sentence leaves
several decisions unresolved: which tests ran, whether they completed,
which artifact was tested, and whether the artifact has since changed.
Similar problems arise when a remembered version is no longer current or
a remembered approval has been revoked. The practical question is
whether the agent can make the next decision using evidence that remains
applicable to the present task state.

The general problem is already well established. Compaction-Friendly
documents the loss of process constraints even when benchmark
performance changes little; CompAct-Bench evaluates downstream
continuation and identifies distorted task state after compaction
(\citeproc{ref-salle2026compactionfriendly}{Salle et al. 2026};
\citeproc{ref-chen2026compactbench}{Chen et al. 2026}). Work on the
Compaction Cliff also studies retention policies that depend on
information type (\citeproc{ref-zerhoudi2026compactioncliff}{Zerhoudi,
Mitrovic, and Granitzer 2026}). Consequently, task success alone is not
a sufficient novelty claim for another study of information loss, and
another example of a forgotten constraint is not the proposed
contribution here.

Nor is retrieving original history a new solution. Raw-RAG demonstrates
the strength of dense retrieval and reranking over untransformed
histories (\citeproc{ref-bao2026rawrag}{Bao et al. 2026}). TierMem
already routes from summaries to original logs when summary evidence is
insufficient (\citeproc{ref-zhu2026tiermem}{Zhu et al. 2026}). HELM
governs memory using provenance, status, freshness, and consistency; GPM
defines source-bound lifecycle and structured-release conditions at a
verified ledger head (\citeproc{ref-liu2026helm}{Liu et al. 2026};
\citeproc{ref-xu2026gpm}{G. Xu 2026}). These results substantially
narrow the research question: \textbf{under native host compaction, does
explicit verification of executable, changing evidence improve an
agent's actual subsequent decisions beyond strong retrieval and
established sufficiency or governance policies?}

We investigate this question through an auxiliary layer, provisionally
called \textbf{Memex verification}. The layer does not require replacing
the host's compactor. It observes the user request, exposed action plans
or response drafts, current workspace state, and authorized source
history. It discovers claims that matter to a pending decision, obtains
original evidence, checks its qualifications, and supplies the result to
the same host. The host remains responsible for generating the
continuation. A verified evidence packet is therefore an intermediate
artifact; whether the host uses it correctly is an empirical endpoint.

Three distinctions guide the design. First, \textbf{relevance is
different from qualification}: a test log can be highly relevant while
describing the wrong code revision. Second, \textbf{local sufficiency is
different from currentness}: an unexpired approval already present in
context can be superseded by a later archive event. Third, \textbf{a
valid packet is different from a correct continuation}: a model can
ignore valid evidence, misstate it, or introduce a new unsupported
assertion. Each distinction must be measured without giving the proposed
method privileged labels or source locations.

The intended contribution is a testable composition in a specific
deployment setting, not a claim to have invented provenance, selective
retrieval, typed memory, or version watermarks. The proposed study has
three commitments:

\begin{enumerate}
\def\labelenumi{\arabic{enumi}.}
\tightlist
\item
  \textbf{An executable verification specification without evaluator
  privileges.} Claims and source locations must be discovered from
  deployment-visible inputs; type-specific checks cover tests, versions,
  and historical authorization.
\item
  \textbf{A precisely bounded reuse condition.} A no-reopen decision
  requires a complete previously verified source view and proof that the
  relevant view and applicability conditions still hold. Missing
  freshness information produces an explicit inability to verify.
\item
  \textbf{A paired estimate of incremental value.} Evaluation separates
  retrieval coverage, semantic qualification, selective acquisition, and
  model consumption, with a strong Raw-RAG comparison and complete
  resource accounting.
\end{enumerate}

None of these empirical commitments has yet been satisfied by the
available project data. If strong retrieval accounts for the observed
benefits, the appropriate outcome is a boundary or negative-result paper
describing when additional verification is unnecessary. The method
should not become more elaborate merely to preserve an anticipated
positive conclusion.

\section{2. Related Work and the Remaining Empirical
Question}\label{related-work-and-the-remaining-empirical-question}

\subsection{2.1. Raw history and selective evidence
acquisition}\label{raw-history-and-selective-evidence-acquisition}

Raw-RAG leaves state interpretation to the reader and uses dense
retrieval plus cross-encoder reranking over source history. Its evidence
analysis emphasizes coverage and the relationships needed to derive
answers (\citeproc{ref-bao2026rawrag}{Bao et al. 2026}). Our main
comparison must therefore preserve raw information, temporal metadata,
retrieval strength, and the actual text delivered to the reader. A weak
keyword baseline cannot establish the proposed increment.

TierMem is a particularly close predecessor: it uses a provenance-linked
summary tier, a runtime sufficiency router, escalation to immutable raw
pages, and optional writeback. Its main benchmark configuration disables
online writeback (\citeproc{ref-zhu2026tiermem}{Zhu et al. 2026}). We do
not claim that asking when to reopen a source is new. The question here
is whether a complete freshness condition and type-specific execution
semantics improve native continuations beyond a sufficiency router. That
difference must be demonstrated experimentally; the title is a research
question, not a priority claim.

MSS-Complement formulates state-conditioned recovery of a sufficient
evidence set, while SERBench separates set recovery from candidate
discovery (\citeproc{ref-feng2026missingcomplement}{Feng et al. 2026}).
We adopt the need to distinguish those stages. Our claimed increment
cannot rest on selecting a small evidence set for a pending decision. It
concerns whether the selected set is still eligible to support that
decision under the current source view and object state.

\subsection{2.2. Governance, provenance, and changing
state}\label{governance-provenance-and-changing-state}

HELM combines hierarchical memory, auditable provenance, validity
signals, and a memory-operation controller. Its paper includes
tool-status and version checks, governance ablations, and a
training-free variant (\citeproc{ref-liu2026helm}{Liu et al. 2026}).
This is direct overlap, not a background system that can be dismissed as
merely a retriever. A faithful available configuration should be
evaluated where feasible; an independently written approximation must be
named as such.

GPM specifies source admission, conflict handling, non-revival after
retraction or deletion, and closure of structured outgoing claims over a
fresh view at one stable verified head. It explicitly separates its
deterministic service results from unrestricted language-model accuracy
(\citeproc{ref-xu2026gpm}{G. Xu 2026}). Our watermark rule is a
correctness precondition of this established kind. The proposed
empirical question additionally includes noisy claim discovery,
heterogeneous original logs, and a free-form host continuation; success
at the structured service boundary does not settle that question.
Conversely, the existence of that additional interface does not
establish that our method is superior.

MemIR separates evidence, retrieval cues, and truth-bearing claims
through a typed representation (\citeproc{ref-jin2026memir}{Jin et al.
2026}). Eywa similarly grounds extracted memories in source evidence and
typed signals (\citeproc{ref-joshi2026eywa}{Joshi 2026}). Agent Zero
Memory includes an intent gate and a reading discipline that limits
citations to evidence actually opened (\citeproc{ref-wu2026agentzero}{Wu
and Zhu 2026}). These precedents rule out claiming typed provenance, an
opened-source requirement, or memory gating alone as new contributions.

MemoRepair withdraws affected descendants before rebuilding and
republishing validated successors
(\citeproc{ref-zhao2026memorepair}{Zhao et al. 2026}). StateFuse makes
contradictions visible and evaluates correction and abstention under
matched policies (\citeproc{ref-volkov2026statefuse}{Volkov, Li, and Luo
2026}). They motivate treating invalidation and unresolved conflicts as
explicit states. We do not introduce a new distributed merge protocol or
general descendant-repair algorithm. The authorization experiments also
remain distinct from actual runtime permission: work on provenance
laundering shows why remembered text must not acquire authority merely
through consolidation (\citeproc{ref-xuj2026provenancelaundering}{J. Xu
et al. 2026}).

\subsection{2.3. Recoverable context and native
compaction}\label{recoverable-context-and-native-compaction}

Folio records linked notes, event manifests, and source records, then
programmatically folds and recovers context through stable handles
(\citeproc{ref-ruan2026folio}{Ruan et al. 2026}). Our target setting is
a host whose native compaction is retained. This is an interface
constraint, not evidence that auxiliary verification outperforms a
recoverable context manager. A direct comparison would require a
separately labeled scaffold-level experiment.

Table 1 makes the intended difference explicit. Its final column states
an unanswered question, not a feature absent from every prior system.

\Needspace{10\baselineskip}

\textbf{Table 1. Closest prior contributions and the evidence this study
must add.}

\begin{longtable}[]{@{}
  >{\raggedright\arraybackslash}p{(\columnwidth - 4\tabcolsep) * \real{0.3333}}
  >{\raggedright\arraybackslash}p{(\columnwidth - 4\tabcolsep) * \real{0.3333}}
  >{\raggedright\arraybackslash}p{(\columnwidth - 4\tabcolsep) * \real{0.3333}}@{}}
\toprule\noalign{}
\begin{minipage}[b]{\linewidth}\raggedright
Prior line
\end{minipage} & \begin{minipage}[b]{\linewidth}\raggedright
Already established design or question
\end{minipage} & \begin{minipage}[b]{\linewidth}\raggedright
Unanswered comparison targeted here
\end{minipage} \\
\midrule\noalign{}
\endhead
\bottomrule\noalign{}
\endlastfoot
Raw-RAG & Strong retrieval over original history & Does verification
improve decisions with the same sources, reader, and budget? \\
TierMem / MSS-Complement & Sufficiency routing and decision-conditioned
evidence acquisition & Does current source eligibility add value beyond
sufficiency alone? \\
HELM / MemIR / Eywa & Epistemic governance, typed provenance, and source
grounding & Does the proposed lightweight integration retain value on
natural native continuations? \\
GPM / StateFuse / MemoRepair & Structured state, conflict and
invalidation contracts & How much correctness survives extraction,
source discovery, and free-form generation? \\
Folio & Recoverable context folding and original-source handles & What
can an auxiliary verifier achieve while retaining native compaction? \\
Compaction-Friendly / CompAct-Bench / Compaction Cliff & Constraint
loss, state distortion, and downstream compaction evaluation & Which
verification mechanisms causally improve the next decision? \\
\end{longtable}

\section{3. Problem Formulation and
Scope}\label{problem-formulation-and-scope}

\subsection{3.1. Observable state and archive
view}\label{observable-state-and-archive-view}

At a decision boundary, let \(x_t\) contain the current request,
observable host context, exposed proposed action or response draft, tool
descriptions, and accessible workspace observations. The method cannot
read hidden evaluator labels, private model reasoning, or undocumented
host state. A native compaction boundary is established through the
host's supported interface and observable completion receipt. A
generated summary or a scripted prompt replacement is not
interchangeable with this condition.

The source archive contains immutable event revisions \(A_h\) through a
verified ingestion head \(h\). Events preserve raw content, source role,
object or session identity where available, original timestamps,
ingestion order, tool execution metadata, and resolvable locations. The
archive includes corrections and withdrawals as events rather than
silently replacing their predecessors. Its guarantee is explicitly
limited to the ingested, authorized view. An external change that has
not been ingested is outside that guarantee.

A trustworthy archive watermark is a runtime-generated commitment to
that view, not an annotation of the expected answer. The conservative
form contains an archive identity and epoch, a monotonically advancing
sequence, a content commitment, and ingestion completeness information.
It is generated over all admitted events. A watermark derived only from
evaluator-labeled relevant records is prohibited. All comparison arms
receive the same source and watermark information.

\subsection{3.2. Candidate claims and decision
obligations}\label{candidate-claims-and-decision-obligations}

A discovered candidate claim is a tuple

\[
c=(\text{type},\text{entity},\text{scope},\text{predicate},
\text{candidate value},\text{as-of time},\text{dependent action}).
\]

Fields may be unresolved. The extractor may propose that a particular
test passed; that proposal is neither evidence nor truth. A version
mentioned by a user can identify a candidate to examine but cannot
itself certify the current version. An ambiguous entity or scope remains
ambiguous until resolved from visible material.

Hidden \textbf{decision obligations} are defined independently by
evaluators before continuations are generated. They identify the
critical questions and actions required for task success. They are used
only for scoring. This separation prevents an extractor from obtaining a
perfect score by omitting difficult claims, and prevents a cautious
answer from winning merely by saying nothing.

\subsection{3.3. Verification outcomes}\label{verification-outcomes}

The verifier returns a resolved state when justified, together with
original evidence and its applicability interval. For a proposed
assertion it can report supported, contradicted, or unresolved. A
contradicted candidate can still yield a useful known answer: an
approval claim can be contradicted by a verified revocation. The host
should not turn every contradiction into generic uncertainty.

Unresolved results include a reason, such as missing source, ambiguous
entity, incomplete event, unavailable current view, unresolved conflict,
or exhausted verification budget. A source that explicitly reports
\texttt{UNKNOWN} is different from a verifier that cannot inspect the
source. Likewise, absence from search results is not proof that a test
did not run. \texttt{NOT\_RUN} requires an authoritative execution
record or a complete execution view establishing non-execution under the
task's semantics.

The unit of practical interest is the resulting task continuation. We
score whether the host selects and executes the required next action and
accurately communicates the relevant state. Merely producing a
structurally valid certificate is not task success.

\section{4. Proposed Method}\label{proposed-method}

This section specifies the \textbf{full proposed deployment protocol}.
Present-tense descriptions define its required behavior; they do not
assert that the current prototype implements every step. Section 5.3
records the narrower implemented boundary, including live binding for
explicitly current-workspace test claims. General object-state binding,
lifecycle projection from natural logs, and independently evaluated
native continuations remain requirements for further work.

\subsection{4.1. Discover claims without evaluator
privileges}\label{discover-claims-without-evaluator-privileges}

The extractor consumes only \(x_t\) and public tool schemas. It
identifies candidate claims from the current request, the exposed action
plan, and any response draft. It returns the supporting request spans,
unresolved fields, and the action whose preconditions depend on each
claim. A frozen prompt or constrained parser may be used, but all
extraction errors count in the end-to-end result.

The production interface must not consume \texttt{case.claims}, correct
source identifiers, correct event indices, hidden required-source lists,
or gold normalized entities. Source locators may appear only after a
real archive operation returns them, or when they were already genuinely
visible to the host. A diagnostic oracle condition may supply correct
claims and sources, but it is reported separately and is never called
the deployable method.

Arms C and D, and any claim-driven ablations, share the same extractor
and visible metadata. The primary faithful Raw-RAG arm B retains the
original user request as its query. Any optional claim-expanded Raw-RAG
condition must be separately named, frozen, and budgeted as an
adaptation; it is not an already implemented baseline or a replacement
for B. The evaluator scores extraction recall against all frozen
obligations, not just against the candidates the method emitted.

\subsection{4.2. Preserve original events and find
sources}\label{preserve-original-events-and-find-sources}

The archive adapter ingests actual tool responses, user corrections,
version observations, and authorized historical records. Derived
summaries and model statements retain their source roles. They cannot
silently become tool execution receipts or user decisions. Ingestion
preserves content hashes and original event boundaries; it does not
require every log to follow a synthetic receipt template.

Retrieval starts from the discovered claim and exposed context. It may
use dense retrieval, lexical retrieval, and reranking, but query
construction cannot use a hidden correct location. Every hit returns a
resolvable locator. The opening layer follows returned pagination and
checks record identity, revision, offsets, content integrity, and
completion. Event position is variable: there is no privileged
\texttt{eventIndex\ =\ 1}. A clipped record is not complete evidence
merely because its visible prefix contains a reassuring word.

Parsing recognizes frozen families of natural stdout, structured tool
responses, test reports, version readbacks, and user corrections.
Unsupported formats remain unresolved. Parsers must derive field values
from original content and trusted tool metadata; a search snippet or an
assistant's paraphrase does not certify execution. Entity normalization
and alias handling are part of the measured method, with their own
failure labels.

\subsection{4.3. Separate positive support from lifecycle
completeness}\label{separate-positive-support-from-lifecycle-completeness}

Finding a positive passage answers whether some source supports a
statement. It does not show that no later event invalidates that
support. In particular, top-k retrieval can return an old grant while
omitting a revocation already present at the same archive head. A stable
head alone does not repair this omission.

A current-state certificate therefore requires both original support and
\textbf{coverage of the applicable lifecycle} through its stated head.
This can be provided by a complete replay of the relevant source lineage
from an independently constructed event index, or by enumerating the
available archive and establishing the required scope. The index must be
built from all admitted events, not from evaluator-provided relevance
labels. Its source spans, parser version, coverage cursor, and
completeness status must be auditable.

Where exact lineage coverage cannot be guaranteed, the full protocol
requires a scan of the permitted event sequence or an unresolved result.
If a scan exceeds the fixed budget, the method records the budget
failure; it must not promote a partial search into a completeness proof.
The formal claim in Section 4.6 is conditional on correct normalization
and complete coverage. Measuring how often those conditions can be
obtained cheaply is part of the research question.

\subsection{4.4. Type-specific
qualification}\label{type-specific-qualification}

The initial method covers three narrowly defined claim families. Table 2
specifies the required semantics. It does not assert a general
natural-language entailment solver.

\Needspace{10\baselineskip}

\textbf{Table 2. Qualification requirements and representative rejection
cases.}

\begin{longtable}[]{@{}
  >{\raggedright\arraybackslash}p{(\columnwidth - 6\tabcolsep) * \real{0.2500}}
  >{\raggedright\arraybackslash}p{(\columnwidth - 6\tabcolsep) * \real{0.2500}}
  >{\raggedright\arraybackslash}p{(\columnwidth - 6\tabcolsep) * \real{0.2500}}
  >{\raggedright\arraybackslash}p{(\columnwidth - 6\tabcolsep) * \real{0.2500}}@{}}
\toprule\noalign{}
\begin{minipage}[b]{\linewidth}\raggedright
Claim family
\end{minipage} & \begin{minipage}[b]{\linewidth}\raggedright
Required positive support
\end{minipage} & \begin{minipage}[b]{\linewidth}\raggedright
State and applicability checks
\end{minipage} & \begin{minipage}[b]{\linewidth}\raggedright
Insufficient evidence
\end{minipage} \\
\midrule\noalign{}
\endhead
\bottomrule\noalign{}
\endlastfoot
Test status & A real execution event, command or suite identity,
completion, and outcome & Match the tested artifact digest, relevant
environment, command semantics, and decision scope; account for later
changes or reruns & Assistant says ``passed''; process only started; no
exit record; ignored subprocess failure; result for an old artifact;
partial output \\
Current version & An authoritative readback for the named object or a
complete ordered update lineage & Resolve source authority and object
identity; respect effective time, supersession, and conflicts &
Candidate version in the question; a similar filename; stale metadata;
two unresolved authoritative values \\
Historical authorization & An attributable decision with principal,
target, scope, conditions, and effective interval & Replay grants,
denials, revocations, expiration, and valid supersession in the ingested
view & Different principal or target; expired grant; later revocation;
ambiguous conditions; quoted text treated as live permission \\
\end{longtable}

For tests, a command's zero exit status only supports the scope
warranted by that command and its wrapper. A shell that ignores a
failing subprocess does not establish that all intended tests passed. A
run against one commit does not certify a modified working tree. A
deployment adapter must record how artifact and environment identifiers
are measured; missing identifiers cannot be replaced by a guessed
version label.

For current versions, the method discovers the observed value rather
than validating a value supplied as the answer. Conflicting observations
are resolved only by frozen authority and lifecycle rules. A later
wall-clock timestamp from a lower-authority source is not a universal
override rule. Ingestion order and event effective time are recorded
separately so late-arriving corrections can be handled without
pretending they arrived earlier.

For authorization, the output is a fact about a specified historical or
simulated state. A legitimate statement that a grant existed at an
earlier time can coexist with a current revocation. All action
consequences in the study are simulated. Neither an archived user
message nor a verification certificate grants permission to execute a
real external action; runtime authorization remains with the host's
current instruction and permission system.

\subsection{4.5. Decide whether to reopen a
source}\label{decide-whether-to-reopen-a-source}

A cached certificate binds the normalized claim, original support
locators and hashes, lifecycle coverage, archive head, object state,
parser and policy versions, and the interval over which its result
remains applicable. It also records whether the support was actually
opened. A fragment copied into compacted context without its verifiable
certificate is not eligible for the fast path.

Before reusing a certificate, the full protocol requires a current
trusted archive watermark and an independent observation of the
applicable current object state. A version supplied in a candidate claim
or a past capture does not, by itself, establish the live object state.
Reuse requires all of the following:

\begin{enumerate}
\def\labelenumi{\arabic{enumi}.}
\tightlist
\item
  The archive identity, epoch, and committed view match the certificate,
  or a complete validated delta has re-established its result.
\item
  The original support and lifecycle coverage remain intact and
  resolvable.
\item
  The relevant artifact, environment, identity mapping, parser, and
  policy versions still match.
\item
  The query scope is unchanged and the decision time lies inside the
  certified applicability interval.
\item
  No unresolved conflict, invalidation, ingestion gap, or inaccessible
  update prevents qualification.
\end{enumerate}

The applicability interval must end at the next known relevant temporal
transition, not merely at a grant's expiration. A future-dated
revocation that is already in the archive can become effective without
another append. The fast path therefore rechecks temporal applicability
even when the watermark is unchanged.

The proposed conservative deployment policy uses a global archive head.
Any appended event invalidates the no-reopen shortcut until its effect
has been checked. Per-entity watermarks may reduce unnecessary
revalidation, but require a proven complete mapping from every event to
its affected entities and policies. They are an optional optimization,
not a hidden assumption. Epoch and content commitments prevent a reset
or rollback from masquerading as an unchanged sequence number.

This is \textbf{selective source reopening}, not zero-I/O verification.
Watermark reads, object observations, delta scans, certificate
validation, and failed attempts all count toward cost. A TTL can impose
an additional refresh policy; it cannot substitute for update detection.

The following pseudocode summarizes the target protocol. The current
prototype implements a bounded form of live observation for
current-workspace test claims; it does not implement this entire
protocol for arbitrary objects and claim families.

\begin{verbatim}
Input: observable task state x, original-source archive A, certificate cache K
C := extract candidate claims from x using the frozen deployment-visible interface
for each candidate c in C:
    read a trusted current archive head h_before and applicable object state
    if a certificate in K has complete coverage and passes all reuse conditions:
        construct a snapshot-bound verification packet from that certificate
    else:
        retrieve candidate sources without gold locations
        open original records and validate their identity, integrity, and completion
        establish complete applicable lifecycle coverage through h_before
        apply the frozen type-specific qualification rules
        construct a supported, contradicted, or unresolved verification packet
    read the archive head again
    if the relevant view changed or completeness cannot be established:
        retry only within the frozen budget; otherwise return unresolved
    record the packet, its evidence, its head, and every incurred operation
deliver packets through the host's supported continuation interface
observe the host's next actions and final claims; revalidate at any current-state commit boundary
\end{verbatim}

New critical claims introduced by a later exposed draft trigger the same
bounded process. If no draft interception is supported, the system
cannot claim complete online control over the host's free-form output.
Such a host may still be evaluated for decision support, with final
claim checking performed by the evaluator and this limitation stated
explicitly.

\subsection{4.6. Conditional reuse
property}\label{conditional-reuse-property}

Let \(R(c,A_h,z,t;p)\) be the deterministic qualification result for a
normalized claim \(c\), complete archive view \(A_h\), applicable object
state \(z\), decision time \(t\), and parser/policy version \(p\).
Suppose a certificate was correctly constructed from the complete
applicable lifecycle at \((h,z,p)\), and it binds an interval \(I\) on
which the result does not change. If a trusted read establishes the same
archive view, \(z\) and \(p\) are unchanged, support integrity holds,
and \(t\in I\), the certified result equals recomputation of \(R\) for
that view. The argument follows directly because all inputs that can
change the specified result are unchanged.

This is a conditional consistency property, not a new general theorem
about memory or a proof of implementation correctness. It assumes
faithful ingestion, correct claim normalization, complete lifecycle
coverage, correct type semantics, and reliable version observations. It
does not prove that a source tells the truth, that unobserved external
state is unchanged, or that the model's final answer follows the
certificate. GPM already makes a closely related stable-head release
discipline explicit (\citeproc{ref-xu2026gpm}{G. Xu 2026}).

A packet states the archive head and as-of interval it certifies. If the
intended action requires current state at commitment, the layer must
revalidate immediately at that supported boundary. Otherwise the claim
is only about the recorded snapshot. An update after a successful head
check is not retroactively covered by that check. This distinction
prevents a packet verified before a long model continuation from being
presented as an indefinite authorization or freshness guarantee.

\subsection{4.7. Deliver evidence and audit
consumption}\label{deliver-evidence-and-audit-consumption}

Packets include the candidate claim, resolved state or reason for
uncertainty, original excerpts, source locators, applicability
information, completeness status, and the certified head. They preserve
contrary evidence when a conflict remains. The host receives these
packets through its supported context or tool interface; the harness
records the delivered payload hash and the host receipt.

We audit five distinct stages: obligation discovery, source discovery,
qualified original opening, actual delivery, and the host's decision or
claim. We do not infer original-source use solely from a correct answer.
Conversely, unavailable internal model telemetry is labeled unobserved
rather than automatically scored as an incorrect decision. The primary
endpoint uses observable actions, final output, and independently
checkable source support.

\section{5. Available Development
Evidence}\label{available-development-evidence}

\subsection{5.1. Earlier synthetic native-host
study}\label{earlier-synthetic-native-host-study}

The versioned project report describes 24 synthetic development
fixtures, each run under a native condition and two archive-assisted
conditions, for 72 runs and 456 native compactions. On its scored
supported conclusions, the reported counts are 24/120 for the native
condition and 120/120 for each archive-assisted condition. Behavioral
code checks pass on all 24 tasks in every condition, and all three
conditions have zero unsupported high-impact assertions in the report's
184 scored opportunities (\citeproc{ref-memex2026pilot}{Memex project
2026}).

These numbers are motivating development observations only. The archive
intervention occurred after code implementation, during a read-only
handoff. They therefore do not demonstrate improved coding success or
reduced harmful-assertion rates. The 120 conclusions and 456 boundaries
are nested observations, not independent sample sizes. The fixtures and
their receipt structures cannot be reused as the proposed confirmatory
population.

The inspected driver consumes predefined \texttt{case{[}"claims"{]}},
constructs a structured source filter, and opens a fixed event position.
It does perform real archive operations; describing it as completely
fabricated retrieval would be inaccurate. Its limitation is that claim
discovery, natural-log parsing, and unrestricted source localization are
not evaluated. The present method removes these conveniences as
deployment inputs. The old study's forced-verification condition must
not be relabeled as Raw-RAG.

\subsection{5.2. Recorded freshness
counterexample}\label{recorded-freshness-counterexample}

A recorded local probe presents an already opened, unexpired
\texttt{granted} event to the old selective gate. Its prepared archive
callbacks expose a later \texttt{revoked} event for the same scope and
object. The recorded gate result is \texttt{supported} with status
\texttt{granted}, \texttt{triggered\ =\ false}, and zero search calls.
The source commit and exact probe record are identified in Appendix C.

This is a deterministic development counterexample, not a prevalence
estimate or a new end-to-end result. It demonstrates a missing
precondition for that gate's no-search branch when the caller has not
independently guaranteed freshness. It does not invalidate the earlier
cohort's reported scores, whose sufficient-evidence controls did not
include that later superseding event. The counterexample must become a
regression test before formal experiments begin.

\subsection{5.3. Current implementation
boundary}\label{current-implementation-boundary}

A separate working verifier is being developed for the P1 study while
preserving the legacy experiment. Its local regression checks establish
bounded implementation behavior, not model-based results or support for
H1/H2. The engineering checks described here are outside the study
denominators and the result tables in Section 7.

The current archive adapter rereads its trusted catalog and computes a
session-wide watermark over all admitted catalog entries, including
entries unrelated to the queried entity. For the relevant normalized
snapshot head, it also returns the sequence, event count, and a
commitment over sorted event-identifier and source-digest pairs. The
verifier checks the full set of opened head events against that count
and commitment and rechecks the view before delivery. Reuse binds the
cached source to the trusted catalog commitment; a cached payload and
its own matching self-reported hash are insufficient. These checks
depend on the trust and completeness of the admitted catalog.

The normalized-source API currently assumes \textbf{complete state
snapshots}, with source sequences ordering those snapshots. It does not
construct a general lifecycle projection from arbitrary natural-language
events. Opening all events in a latest append batch would not establish
full history unless that batch satisfies the complete-snapshot contract.
The narrow Node/TAP adapter parses supported native command events,
execution completion, exit status, and a nonzero executed-test count,
and checks equal artifact hashes captured before and after execution.
The capture metadata must come from a trusted host observer.

A live-state component now observes an explicit bounded list of public
artifact files and lockfiles, the observing Node runtime, a declared
relative execution directory, and declared public environment values. It
hashes file contents and observed directory membership, rejects
noncanonical paths and observed symlinks, and records identity and
metadata change guards. An execution source is admitted as live-bound
only when successful before/after observations agree and the normalized
source already carries that observed version. The claim-binding
interface explicitly requires a current-workspace target, retains the
model-proposed version, and rejects explicit historical queries rather
than silently changing their time scope.

The live wrapper observes the workspace before and after every archive
view read and includes its binding summary in the view. It validates a
new committed head through real, bounded original-source search and
opening before permitting cached delivery. Consequently, an observed
file change invalidates old cached support even if the archive catalog
has not changed; rebinding a claim to the new version does not make an
old test receipt applicable. Fifteen local tests pass for this
component, covering content changes with restored mtime,
directory-member additions and deletions, lockfile changes, unavailable
files, path checks, observed races, execution binding, and
cached-delivery invalidation through actual L2 storage. These are finite
engineering checks. Repeated observations are not an atomic snapshot or
a substitute for operating-system isolation, and the host must align the
actual execution runtime and directory with the observer.

An observable-input boundary and output validation for claim extraction
are implemented, but they do not establish a model's extraction accuracy
on the proposed natural tasks. A separate controller implements a local
content freeze, bounded dispatch, retained first-attempt records, an
evaluator-only join, and component-level usage ledgers for later
analysis. Its contract checks use temporary synthetic fixtures and
controlled runner boundaries and have been revised following review. A
local freeze is not preregistration, a declared runner boundary is not
filesystem isolation, and a ledger is not proof that telemetry or
independent labels are complete. No controller test count or measured
arm effect is used as a study result here.

\textbf{Online integration status.} One new macOS native-host
development attempt started a generation turn and successfully called
the real MCP workspace-inspection tool. Its test-tool call was rejected
because approval was required but the unattended runner used a policy
that could not grant it. That attempt executed no test command and
reached no compaction, claim extraction, or verification decision. The
blocked attempt is retained; it is not a completed method-arm run or
evidence of end-to-end integration. Earlier native-host smoke receipts
are separate records and cannot fill these missing stages.

A separate extraction-only component attempt ran one real model
generation without an MCP server, test execution, compaction, or
verification decision. Its output passed JSON and schema validation but
omitted the required test-result candidate, listing that obligation only
as unresolved. It emitted an authorization candidate and two version
candidates and failed the component's test-obligation check. A general
prompt revision then clarified that decision preconditions are
candidates requiring verification even when no result is available, and
that the command identifies the test entity rather than the artifact
version. One retest on the same synthetic development input emitted one
candidate from each of the three claim families, including the required
test candidate with its command entity and proposed artifact digest. It
passed the narrow component check while explicitly leaving coverage
unestablished. Both attempts are retained; the retest does not replace
the initial semantic failure or provide a natural-task accuracy
estimate.

Across these host-integration and extraction-only attempts, three model
generation turns started, with zero test executions, native compactions,
verification decisions, and formal-study tasks. Provider charges and
standardized cost remain unmeasured (\texttt{null}), not zero. General
semantic normalization, lifecycle completeness, other object-state
bindings, and independently evaluated native continuations remain
acceptance requirements. Appendix C identifies the retained attempt and
component records.

\subsection{5.4. Source acquisition and environment
qualification}\label{source-acquisition-and-environment-qualification}

The source-acquisition record contains 16 distinct issue/PR candidates
from eight projects in a fixed revision of SWE-bench Verified. It
records upstream identities, base and fix commits, patch-field digests,
license evidence, and official image-manifest digests. It currently
contains \textbf{zero accepted or frozen P1 workflows} and no
independent P1 double annotation. These are convenience-sampled sources
for workflow construction, not 16 completed pilot tasks. Dataset
curation does not supply the study's temporal obligations,
evidence-equivalence rules, or post-compaction hidden behavior. Appendix
C identifies the acquisition and reproduction records.

Two candidates have completed local program reproduction using the fixed
dataset's official test and production patches. For Flask PR 5014, the
base with the test patch produced 1 failed and 59 passed tests; applying
the production patch produced 60 passed. For Requests PR 6028, the
corresponding counts were 2 failed, 201 passed, and 11 skipped, followed
by 203 passed and 11 skipped. These are four pytest invocations across
two repository instances, restricted to the respective target test
files. The counts were recomputed from retained JUnit reports, and the
same test-file contents were used before and after the production patch.
Skips are not passes. The records also preserve patch checks, dependency
versions, phase hashes, command arguments, exit codes, and
environment-preparation failures.

This exercise qualifies sources and a documented local Python
environment. It ran no language model, native compaction, retrieval arm,
or P1 hidden-behavior evaluation, and it supplies no arm success rate,
pilot estimate, or cost endpoint. The local environments are not
asserted to be bitwise-equivalent to official containers. Filesystem and
selected process namespaces were isolated, but the final runs shared the
host network namespace after a full network-isolation attempt was
refused. The evaluator clones have encountered gold tests and contain
future history; they must not become agent-visible workspaces without a
separate audited public export.

Source reconstruction also has material pitfalls. In 11 of the 16
candidates, the dataset base is not the upstream fix's first parent. A
full base-to-merge diff could therefore add unrelated intervening
changes; reconstruction must use the pinned dataset patches and verify
their application. Official test labels also require inspection: one
pytest candidate's pass-to-pass list contains a progress-display token,
while the Requests parser truncates whitespace-bearing identifiers and
normalizes absolute-path parameters to basenames. Multiple actual
parameterized tests can map to one official label. The reproduction
preserves the initial literal-identifier mismatches and the subsequent
application of the pinned official parser, and reports actual JUnit test
counts rather than label counts. Neither unit-test counts nor repeated
phases increase the study's independent task sample size.

\section{6. Planned Evaluation}\label{planned-evaluation}

\subsection{6.1. Tasks and native
continuations}\label{tasks-and-native-continuations}

Tasks will be assembled from licensed, independent repository workflows.
A representative workflow includes an observed failure, a modification,
an evidence or requirement change, native compaction interleaved with
real tool use, and a subsequent decision that depends on prior evidence.
Examples include deciding whether to rerun tests after a code change,
selecting the correct artifact version, or satisfying a simulated
approval condition after a correction. Source forms, event locations,
tool names, and distractor fields will vary.

At least one genuine state change is required in the main task design.
Repeating compaction after all implementation work is complete is a
stress test, not a substitute for an evolving workflow. We will report
natural boundary timing separately from prescribed 3-, 6-, and
10-compaction stress strata. All authorization consequences remain
isolated simulations.

Tasks are split by repository, source event, and workflow family.
Renaming variables in one template does not create independent tasks.
The old 24 fixtures are development-only. The plan starts with 8
implementation tasks and 16 new pilot tasks; the confirmatory sample
size is determined from pilot-based design analysis and frozen before
unsealing outcomes. A separate 40-task external-validation set is
planned, subject to demonstrated independence and available resources.

Each paired continuation begins from the same observable prefix and
environment snapshot where the host provides a supported branch or
restore mechanism. No private host state is modified to manufacture
equivalence. If exact branching is unavailable, a documented replay
interface may establish frozen-trajectory continuations, which are
labeled accordingly and supplemented with natural runs from the
beginning. A host for which comparable starting states cannot be
established is excluded from the main causal contrast.

\subsection{6.2. Comparison arms}\label{comparison-arms}

\Needspace{10\baselineskip}

\textbf{Table 3. Planned comparison arms. No arm has a new empirical
score in this draft.}

\begin{longtable}[]{@{}
  >{\raggedright\arraybackslash}p{(\columnwidth - 4\tabcolsep) * \real{0.3333}}
  >{\raggedright\arraybackslash}p{(\columnwidth - 4\tabcolsep) * \real{0.3333}}
  >{\raggedright\arraybackslash}p{(\columnwidth - 4\tabcolsep) * \real{0.3333}}@{}}
\toprule\noalign{}
\begin{minipage}[b]{\linewidth}\raggedright
Arm
\end{minipage} & \begin{minipage}[b]{\linewidth}\raggedright
Configuration
\end{minipage} & \begin{minipage}[b]{\linewidth}\raggedright
Role
\end{minipage} \\
\midrule\noalign{}
\endhead
\bottomrule\noalign{}
\endlastfoot
A & Native host with its ordinary compaction & Contextual reference for
default behavior \\
B & Strong Raw-RAG over the same original archive, with a frozen
evidence-budget adaptation & Primary comparative baseline \\
C & All discovered claims use the same source-acquisition and
typed-verification pipeline & Always-verify quality and cost control \\
D & Full method, including certified selective reuse & Proposed
deployable method \\
E & Ordinary archive search and open & Small mechanistic control for
archive access \\
F & Shared candidate/source packets with typed state qualification
removed & Isolate the incremental effect of qualification \\
G & Correct claims and required originals supplied together & Combined
diagnostic oracle, never included in deployable results \\
H & Faithfully reproduced available governance configuration &
Close-method comparison; otherwise explicitly a self-implemented
contract-inspired control \\
T & Faithful TierMem-style sufficiency route, if reproducible; otherwise
a clearly named independent sufficiency-only control & Added neighboring
comparison motivated by the literature review \\
\end{longtable}

Arms C and D share extraction, parsers, retrieval, source budget, and
qualification rules. D differs in whether a valid certificate permits
reuse. F uses the same fixed candidate material as its paired qualified
condition, so different evidence content cannot be mistaken for a rule
benefit. All relevant baselines can see the same ordinary source
metadata, time information, and archive head; none receives evaluator
truth.

The Raw-RAG reference pipeline uses 256-token windows with 32-token
overlap, dense top-128 retrieval, and cross-encoder selection of 15
windows. Its selected windows may expand to full turns, so 15 selected
windows must not be described as a fixed 3,840-token reader input
(\citeproc{ref-bao2026rawrag}{Bao et al. 2026}). We freeze the official
code revision, tokenizer, encoders, reranker, query construction,
effective configuration, expansion and ordering behavior, prompt, and
retry policy. Original event timestamps and provenance are preserved;
the adapter may not invent historical timestamps for undated records.

The primary B configuration uses a prespecified assembly policy after
source expansion to match the study's evidence budget. A separate
faithful official presentation configuration is run if it differs
materially from B. If these configurations are byte-identical under the
same prompts and sources, their runs may be reused; otherwise they
require distinct runs and resource entries. Tuning occurs only on
development data.

H requires an artifact and a conformance record. GPM's inspected v2
paper states that its complete reproducibility artifact had not yet been
publicly uploaded, so a faithful GPM reproduction cannot be promised on
that basis (\citeproc{ref-xu2026gpm}{G. Xu 2026}). A conservative
implementation of its stated ideas is not named ``GPM'' in results.
HELM's available training-free or governance-only configuration is a
candidate for reproducibility assessment, not a baseline already
reproduced here.

Arm T is an addition to the initial planning document. Its exact
implementation, subset or full-set coverage, and extra cost must be
frozen before evaluation. It may not be added only after observing a
favorable result. If a close method cannot be run, the paper must retain
the comparison limitation instead of presenting Raw-RAG as the only
relevant predecessor.

\subsection{6.3. Primary endpoint}\label{primary-endpoint}

For task \(i\), arm \(a\), and stochastic repetition \(r\), let
\(Y_{iar}=1\) when the continuation satisfies all frozen critical
decision obligations, passes the relevant hidden behavioral checks,
makes no incorrect high-impact affirmative assertion, and supports its
critical affirmative claims with evidence applicable to the current
object and state. Otherwise \(Y_{iar}=0\).

An unnecessary abstention on an answerable obligation is a failure. When
an obligation is genuinely unanswerable from the permitted view, an
appropriately explained abstention may satisfy it. Missing a required
decision is a failure even if every sentence the model does emit is
technically true. Citations may use any equivalent resolvable locator; a
particular receipt spelling is not required. A correct unsupported guess
does not satisfy an evidence-support obligation.

This composite endpoint is named \textbf{evidence-supported decision
success}. Its components are also reported separately, because the same
aggregate failure can arise from behavioral failure, omitted
obligations, wrong values, unsuitable sources, or over-abstention.
Scoring rules distinguish correct state resolution from whether a source
was actually retrieved and delivered.

\subsection{6.4. Hypotheses and
estimands}\label{hypotheses-and-estimands}

\textbf{H1: incremental effectiveness.} The full method improves mean
task-level evidence-supported decision success over B. The primary
estimand is

\[
\Delta_{DB}=\frac{1}{N}\sum_{i=1}^{N}
\left(\overline{Y}_{iD}-\overline{Y}_{iB}\right),
\qquad
\overline{Y}_{ia}=\frac{1}{R}\sum_{r=1}^{R}Y_{iar}.
\]

The planned design uses two randomized repetitions per task and arm, but
the independent unit remains the task. Repetitions, claims, and
compaction boundaries do not increase \(N\). We report the paired effect
and its 95\% confidence interval and use a prespecified two-sided
\(\alpha=0.05\) test.

\textbf{H2: selective cost reduction with quality noninferiority.}
Following support for H1, compare D with C. The proposed absolute
noninferiority margin is three percentage points, which requires a
substantive justification and a frozen value before the pilot informs
the confirmatory design. The quality condition requires the lower
endpoint of the two-sided 95\% interval for \(\Delta_{DC}\) to be
strictly greater than \(-0.03\). The cost condition requires the upper
endpoint of the 95\% interval for mean paired standardized total
processing cost \(C_D-C_C\) to be strictly below zero. Both conditions
must hold. Noninferiority is not equality.

\textbf{H3: external validity.} Estimate effect sizes on independent
workflows under a different model and a different host configuration.
These are separately reported external-validity analyses; changing both
host and model is not a pure model intervention. Small favorable
subgroups do not replace the primary comparison.

Repository or source-event clustering is retained. The primary interval
procedure resamples repositories and then tasks, preserving paired arm
outcomes and repetitions within each task. A task-paired sensitivity
analysis is also reported. When the number of repositories is too small
for stable cluster inference, that limitation must be disclosed and the
design expanded or the claim narrowed before confirmation.

\subsection{6.5. Cost, latency, and
reliability}\label{cost-latency-and-reliability}

H2 has one primary cost endpoint: mean standardized total processing
cost per scheduled task. Before confirmatory execution, the protocol
freezes reference rates for uncached input, cache reads and writes,
output generation, extraction, embeddings, reranking, local compute, and
storage. One-time indexing and shared prefix costs follow a
preregistered deployment-based amortization rule. Physical research
expenditure is recorded separately to avoid double-counting a shared
prefix that ran once in the study.

Watermark reads and lifecycle scans belong in this cost. So do failed
calls, retained infrastructure failures, and allowed reruns. Reducing
search calls does not by itself establish lower cost or latency. We
report token components, model and tool calls, source bytes read,
compute time, p50/p95 end-to-end latency, and actual charges where
available. If credible reference rates cannot be frozen, H2 cost is
descriptive; a different primary cost metric is not selected after
results are known.

All arms have frozen token, call, and wall-clock limits. Format errors
and ordinary model failures count as failures. Infrastructure
interruptions, timeouts, and transmission errors remain in the record
under the preregistered intention-to-treat rule, with a separate
valid-run sensitivity analysis. Permitted reruns are governed by a fixed
infrastructure-error policy; unfavorable model outcomes cannot be rerun
selectively. A detected model-service version change suspends the
affected block and creates a separately documented stratum.

\subsection{6.6. Mechanistic analyses}\label{mechanistic-analyses}

The following analyses are specified before confirmation:

\begin{itemize}
\tightlist
\item
  \textbf{Fixed evidence.} Give the same original content and metadata
  to an unqualified reader and the typed verifier. Include correct
  values with wrong sources, stale versions, incomplete tests, and
  support accompanied by counterevidence. Measure value accuracy, source
  entailment, state eligibility, and decision success separately.
\item
  \textbf{Freshness interventions.} Hold the current context fixed while
  adding no change, revocation, supersession, conflict, dropped
  notification, out-of-order arrival, and restart scenarios. Compare the
  old local gate, TTL-only refresh, always-verify, and the certified
  method. Include a revocation already inside the stable head but
  missing from top-k results, and a future-effective event that
  activates without an append.
\item
  \textbf{Discovery and delivery.} Score candidate recall against frozen
  obligations and report source-search, opening, completeness, delivery,
  and model-use failures with explicit denominators. The combined oracle
  G provides a diagnostic ceiling when both discovery and source
  acquisition are corrected.
\item
  \textbf{Compaction and change controls.} On 16 prespecified tasks,
  compare B and D under one variant that retains state changes without
  compaction and another that retains compaction without subsequent
  changes. These are two single-factor variants, not a complete
  factorial experiment; no general interaction estimate is claimed.
\end{itemize}

The main mechanism subset is planned to contain 24 prespecified
confirmatory tasks. It is nested within the primary sample and does not
add independent tasks. Every later exploratory analysis is labeled and
kept separate from prespecified tests.

\subsection{6.7. Sample size and
freezing}\label{sample-size-and-freezing}

The 16-task pilot estimates paired discordance, task and repository
variance, completion rates, and resource distributions. Design
simulations examine H1 effects of 5, 8, and 12 percentage points. A
separate simulation evaluates the joint H2 criterion using the
three-point quality margin and the observed quality-cost dependence. The
confirmatory sample must be adequate for the intended claims, or H2 must
be declared estimation-only before unsealing formal outcomes.

An illustrative \(N=120\) is a resource example, not a power result.
Additional neighboring baselines and official Raw-RAG configurations
increase the run count; they do not increase independent \(N\). If
resources are insufficient, the study narrows its claims or waits for a
sufficiently powered design. It does not lower the quality threshold
after seeing results.

Before confirmation, freeze task definitions, source snapshots,
deployment-visible fields, labels, arms, code, prompts, model and host
versions, configuration, randomization, budgets, exclusion and retry
rules, score computation, and statistical scripts. Record hashes and
have an independent evaluation role hold the labels. Critical states
receive independent double annotation; all between-arm disagreements and
a random 20\% of agreements are reviewed. Model judges may assist triage
but do not replace executable behavior checks or source qualification.

\section{7. Results - Reserved for Future
Measurements}\label{results---reserved-for-future-measurements}

\textbf{No confirmatory or paired method-effect results are available.
\texttt{PENDING} means not measured; it does not mean zero, failure, or
a missing value that can be inferred from the earlier pilot.} The tables
below are analysis templates, not reported findings.

\Needspace{10\baselineskip}

\textbf{Table 4a. Main paired decision and behavior outcomes.}

\begin{longtable}[]{@{}
  >{\raggedright\arraybackslash}p{(\columnwidth - 8\tabcolsep) * \real{0.1875}}
  >{\raggedleft\arraybackslash}p{(\columnwidth - 8\tabcolsep) * \real{0.2500}}
  >{\raggedright\arraybackslash}p{(\columnwidth - 8\tabcolsep) * \real{0.1875}}
  >{\raggedright\arraybackslash}p{(\columnwidth - 8\tabcolsep) * \real{0.1875}}
  >{\raggedright\arraybackslash}p{(\columnwidth - 8\tabcolsep) * \real{0.1875}}@{}}
\toprule\noalign{}
\begin{minipage}[b]{\linewidth}\raggedright
Arm
\end{minipage} & \begin{minipage}[b]{\linewidth}\raggedleft
Independent tasks
\end{minipage} & \begin{minipage}[b]{\linewidth}\raggedright
Evidence-supported decision success
\end{minipage} & \begin{minipage}[b]{\linewidth}\raggedright
Behavior success
\end{minipage} & \begin{minipage}[b]{\linewidth}\raggedright
Incorrect high-impact affirmatives
\end{minipage} \\
\midrule\noalign{}
\endhead
\bottomrule\noalign{}
\endlastfoot
A Native & PENDING & PENDING & PENDING & PENDING \\
B Raw-RAG & PENDING & PENDING & PENDING & PENDING \\
C Always-verify & PENDING & PENDING & PENDING & PENDING \\
D Selective verification & PENDING & PENDING & PENDING & PENDING \\
H Governance comparison & PENDING & PENDING & PENDING & PENDING \\
T Sufficiency comparison & PENDING & PENDING & PENDING & PENDING \\
\end{longtable}

\Needspace{10\baselineskip}

\textbf{Table 4b. Abstention and total-resource outcomes for the same
scheduled tasks.}

\begin{longtable}[]{@{}
  >{\raggedright\arraybackslash}p{(\columnwidth - 6\tabcolsep) * \real{0.2500}}
  >{\raggedright\arraybackslash}p{(\columnwidth - 6\tabcolsep) * \real{0.2500}}
  >{\raggedright\arraybackslash}p{(\columnwidth - 6\tabcolsep) * \real{0.2500}}
  >{\raggedright\arraybackslash}p{(\columnwidth - 6\tabcolsep) * \real{0.2500}}@{}}
\toprule\noalign{}
\begin{minipage}[b]{\linewidth}\raggedright
Arm
\end{minipage} & \begin{minipage}[b]{\linewidth}\raggedright
Unnecessary abstention
\end{minipage} & \begin{minipage}[b]{\linewidth}\raggedright
Mean total standardized cost
\end{minipage} & \begin{minipage}[b]{\linewidth}\raggedright
p95 latency
\end{minipage} \\
\midrule\noalign{}
\endhead
\bottomrule\noalign{}
\endlastfoot
A Native & PENDING & PENDING & PENDING \\
B Raw-RAG & PENDING & PENDING & PENDING \\
C Always-verify & PENDING & PENDING & PENDING \\
D Selective verification & PENDING & PENDING & PENDING \\
H Governance comparison & PENDING & PENDING & PENDING \\
T Sufficiency comparison & PENDING & PENDING & PENDING \\
\end{longtable}

\Needspace{10\baselineskip}

\textbf{Table 5. Prespecified inferential decisions.}

\begin{longtable}[]{@{}
  >{\raggedright\arraybackslash}p{(\columnwidth - 4\tabcolsep) * \real{0.3333}}
  >{\raggedright\arraybackslash}p{(\columnwidth - 4\tabcolsep) * \real{0.3333}}
  >{\raggedright\arraybackslash}p{(\columnwidth - 4\tabcolsep) * \real{0.3333}}@{}}
\toprule\noalign{}
\begin{minipage}[b]{\linewidth}\raggedright
Claim
\end{minipage} & \begin{minipage}[b]{\linewidth}\raggedright
Estimate and 95\% interval
\end{minipage} & \begin{minipage}[b]{\linewidth}\raggedright
Required interpretation
\end{minipage} \\
\midrule\noalign{}
\endhead
\bottomrule\noalign{}
\endlastfoot
H1: D minus B success & PENDING & State the effect and uncertainty; do
not substitute D minus A \\
H2 quality: D minus C & PENDING & Noninferiority only if the lower
endpoint exceeds the frozen margin \\
H2 cost: D minus C & PENDING & Cost reduction only if the upper endpoint
is below zero \\
H2 joint statement & PENDING & Assert only if both conditions hold under
the frozen hierarchy \\
External host/model effects & PENDING & Report every planned
configuration and its limitations \\
\end{longtable}

\Needspace{10\baselineskip}

\textbf{Table 6. Mechanisms and failure stages.}

\begin{longtable}[]{@{}
  >{\raggedright\arraybackslash}p{(\columnwidth - 4\tabcolsep) * \real{0.3333}}
  >{\raggedright\arraybackslash}p{(\columnwidth - 4\tabcolsep) * \real{0.3333}}
  >{\raggedright\arraybackslash}p{(\columnwidth - 4\tabcolsep) * \real{0.3333}}@{}}
\toprule\noalign{}
\begin{minipage}[b]{\linewidth}\raggedright
Mechanism or stage
\end{minipage} & \begin{minipage}[b]{\linewidth}\raggedright
Required denominator
\end{minipage} & \begin{minipage}[b]{\linewidth}\raggedright
Measurement
\end{minipage} \\
\midrule\noalign{}
\endhead
\bottomrule\noalign{}
\endlastfoot
Missed candidate claims & All independently annotated obligations &
PENDING \\
Incorrect candidate claims & All emitted candidates, with independent
adjudication & PENDING \\
Source discovery failure & Obligations requiring archive evidence &
PENDING \\
Incomplete lifecycle qualification & Cases requiring current-state
closure & PENDING \\
Incorrect certified reuse & All reuse decisions; additionally all update
opportunities & PENDING \\
Failure to deliver evidence & Packets intended for host delivery &
PENDING \\
Incorrect model use & Obligations with qualified evidence actually
delivered & PENDING \\
Fixed-evidence rule effect & Paired tasks in the frozen mechanism subset
& PENDING \\
\end{longtable}

The analysis order is fixed: primary D/B effect, selective D/C quality
and cost, mechanisms, and external validity. Descriptive legacy fixture
results remain in Section 5 rather than being mixed into Table 4. If the
native arm again has no unsupported high-impact assertions, the paper
will report that observation and will not claim a relative reduction in
that rate.

\section{8. Discussion: What Different Outcomes Would
Mean}\label{discussion-what-different-outcomes-would-mean}

A positive H1 result would support the usefulness of this particular
verifier under the measured tasks, sources, reader, and budgets. The
fixed-evidence comparison must then determine whether the increment
comes from qualification itself or from an improved evidence set. If the
apparent gain disappears when sources and metadata are identical, the
main result is an acquisition effect and must be described accordingly.

A positive joint H2 result would show that certified reuse reduces the
frozen total cost measure within the declared quality margin. It would
not imply lower cost for every workload, every provider tariff, or
frequent-update archives. A global watermark may be inexpensive to read
yet cause near-universal revalidation when unrelated events arrive. A
more selective watermark can improve efficiency only if its dependency
coverage is independently justified.

If D and B perform similarly with a narrow interval, the study can
identify regimes where capable readers already reconstruct the required
state from strong raw evidence. If the interval is wide, the result
remains uncertain rather than demonstrating equivalence. If D is worse,
the analysis should examine extraction omissions, overly strict
qualification, inaccessible metadata, unnecessary abstention, and the
burden of evidence packets on the reader.

Even a correct verifier may have limited task value when evidence is
delivered too late or the host cannot intercept the relevant action. The
distinction between packet correctness and model behavior is therefore
central. Conversely, a favorable model answer does not repair a verifier
that certifies an old grant despite an ingested revocation. Average
benchmark improvement cannot excuse a known violation of the stated
reuse precondition.

This study also does not absorb the broader research agendas of a
cross-system diagnostic benchmark or a text/KV consistency system. Its
diagnostic measurements exist to test the proposed intervention.
Persistent cache design, distributed update consistency, GPU
performance, and general cross-host rankings require their own
independently scoped evidence.

\section{9. Limitations}\label{limitations}

The present document contains a method specification and a planned
study, not completed empirical support. Its model-based development
evidence remains small, synthetic, and structurally convenient. New
live-state regression checks and two real-repository program
reproductions establish engineering and environment properties only. The
source pool has no accepted frozen P1 workflow or completed independent
pilot. These records provide no estimate of the proposed method's
generalization, incremental performance over Raw-RAG, or total cost
advantage.

The proposed semantics cover three claim families and depend on correct
identification of objects, source roles, time, and execution boundaries.
Natural logs can be ambiguous, and source statements can be false.
Integrity and provenance show what was recorded; they do not establish
world truth. Complete lineage may be expensive or impossible to obtain
under real access restrictions.

Freshness is relative to a clearly identified ingested view. A source
outside that view may have changed. Clock error, missing ingestion
events, ambiguous supersession, unsupported formats, parser mistakes,
and incorrect entity resolution can invalidate the preconditions. The
conditional reuse argument is not an assurance that these engineering
components are correct.

The implemented live binding covers explicitly declared public files and
current-workspace test claims. It does not infer all dependencies,
observe arbitrary external services, or retarget historical queries. Two
matching observations do not prove that no intervening change occurred,
and standard-library path checks do not establish a hostile-filesystem
sandbox. The trusted controller must bind actual execution to the
observed runtime and directory. Retrieval can also fail before
qualification: a development probe with a full executable path in the
command entity returned no full-query hits even though shorter queries
found the source. Both the wrapped and unwrapped verifiers abstained.
This is an unresolved retrieval boundary, not evidence that the missing
source is false or absent.

The candidate pool comes from an existing public benchmark, with unknown
pretraining exposure and project-level dependence. Source and license
records do not establish redistribution rights for every issue body or
dependency. Upstream patch ancestry and parser-label quirks require
explicit reconstruction checks; benchmark labels cannot be imported as
P1 truth. Any candidate used to tune extraction, triggering, or budgets
must be assigned to development and replaced in an independent pilot.

Native hosts differ in observable compaction, branching, output
interception, and telemetry. A fixed-trajectory continuation may not
represent an agent's entire natural trajectory. A host with only
post-hoc output inspection cannot support claims of enforced free-form
release correctness. Model and service versions may also change during
data collection.

Matching an evidence budget does not make every resource equal.
Preprocessing, verification scans, embeddings, and the cost of repeated
context can dominate. The paper must show actual resource components and
must not turn a smaller number of search calls into an unmeasured
efficiency claim. A three-point noninferiority margin is a proposed
design choice that still requires a defensible rationale and adequate
sample size.

Finally, related work already covers many components of the proposal. A
combination of those components is not automatically a publishable
contribution. The study earns its methodological claim only through a
clear interface, independently checkable correctness conditions, and
persuasive evidence of a useful difference or a well-characterized
boundary.

\section{10. Reproducibility and Research
Practice}\label{reproducibility-and-research-practice}

The planned artifact contains source and dependency locks, task and
license manifests, event snapshots, normalization rules, exact prompts,
native-host receipts, arm configurations, randomization, complete
per-task records, failure logs, scoring code, and statistical scripts.
It should support result recomputation without model calls and
distinguish that from full model-based reproduction. Artifact
identifiers, source hashes, original offsets, and delivered-payload
hashes connect reported outcomes to their evidence.

Data collection records usage and redistribution permissions for
repository code, issues, documents, and logs separately. Public
readability is not treated as a blanket redistribution license. Private
materials require research authorization. Release preparation removes
credentials, personal identifiers, and machine-private paths while
retaining an auditable mapping from redacted records to checksums where
appropriate.

Author contributions and AI assistance must reflect the actual workflow.
This draft was prepared with AI assistance in literature inspection,
method articulation, and English drafting. It is not accurate to
describe that assistance as language polishing alone. Human authors must
independently verify references, implementation, data provenance,
analyses, and final claims, and approve any submission. No author
identities or institutional affiliations have been fabricated here.

\section{11. Conclusion}\label{conclusion}

This working draft asks whether source verification can improve native
agent continuations after compaction when strong raw-history retrieval
is already available. The proposed layer discovers claims from visible
task state, checks type-specific execution evidence, and permits reuse
only under an explicit complete-view freshness condition. The open
empirical question is whether those checks improve decisions enough to
justify their own errors and cost. The current development records
establish neither superiority nor efficiency. The next necessary
evidence is independently validated host integration and source
qualification, accepted workflows with separated public inputs and
hidden labels, and the paired evaluation specified above.

\clearpage

\section{References}\label{references}

\phantomsection\label{refs}
\begin{CSLReferences}{1}{0}
\bibitem[\citeproctext]{ref-bao2026rawrag}
Bao, Zhengye, Shengtao Zhang, Haokai Zhang, Shiqi Lou, Muning Wen, and
Minnan Luo. 2026. {`Is Complex Agentic Memory Always Needed? Naive {RAG}
Is Stronger Than We Thought'}.
\url{https://doi.org/10.48550/arXiv.2610.30755}.

\bibitem[\citeproctext]{ref-chen2026compactbench}
Chen, Jiayi, Lizhe Chen, Yujia Chen, Shihan Dou, Ming Zhang, Chendi Ge,
Shiguang Ni, et al. 2026. {`{CompAct-Bench}: A Benchmark for
Working-Context Compaction in Long-Horizon Agent Tasks'}.
\url{https://doi.org/10.48550/arXiv.2611.25673}.

\bibitem[\citeproctext]{ref-feng2026missingcomplement}
Feng, Zhexi, Ruiyi Zhang, Yongbo Yang, and Pengtao Xie. 2026. {`The
Missing Complement: State-Conditioned Minimal Sufficient Evidence for
Coding Agents'}. \url{https://doi.org/10.48550/arXiv.2609.20050}.

\bibitem[\citeproctext]{ref-jin2026memir}
Jin, Zhengda, Bingbing Wang, Jing Li, Ruifeng Xu, and Min Zhang. 2026.
{`Mitigating Provenance-Role Collapse in Long-Term Agents via Typed
Memory Representation'}.
\url{https://doi.org/10.48550/arXiv.2605.25869}.

\bibitem[\citeproctext]{ref-joshi2026eywa}
Joshi, Resham. 2026. {`{Eywa}: Provenance-Grounded Long-Term Memory for
{AI} Agents'}. \url{https://doi.org/10.48550/arXiv.2605.30771}.

\bibitem[\citeproctext]{ref-liu2026helm}
Liu, Jialong, Yuling Shi, Zuchao Li, and Ping Wang. 2026. {`{HELM}:
Steering Long-Horizon Agents with Learned Hierarchical Memory and
Epistemic Governance'}. In \emph{Proceedings of the 2026 Conference on
Empirical Methods in Natural Language Processing}, 16728--53. Budapest,
Hungary: Association for Computational Linguistics.
\url{https://doi.org/10.18653/2026.emnlp-main.721}.

\bibitem[\citeproctext]{ref-memex2026pilot}
Memex project. 2026. {`Memex Native Evidence-Recovery Development
Report'}. Versioned project technical report.
\url{https://github.com/QiuXu-Tech/memex/blob/49e772bbf60532682771874b7ec96958f4a179e3/docs/research/evidence-recovery-2026-11-26/RESULTS.md}.

\bibitem[\citeproctext]{ref-ruan2026folio}
Ruan, Yubo, Tao Peng, Jiao Zhang, Chuanguo Ma, Lei Xu, Han Zhou, Wenwu
Wang, et al. 2026. {`{Folio}: Recoverable Multi-Resolution Context
Folding for Long-Horizon Language Agents'}.
\url{https://doi.org/10.48550/arXiv.2611.16627}.

\bibitem[\citeproctext]{ref-salle2026compactionfriendly}
Salle, Jeanne, Hardik Bhatnagar, Matthias Bethge, Sahar Abdelnabi, and
Maksym Andriushchenko. 2026. {`Benchmarks Are Surprisingly
Compaction-Friendly'}. \url{https://doi.org/10.48550/arXiv.2610.31778}.

\bibitem[\citeproctext]{ref-volkov2026statefuse}
Volkov, Sergey, Yang Li, and Ye Luo. 2026. {`{StateFuse}: Deterministic
Conflict-Preserving Memory for Multi-Agent Systems'}.
\url{https://doi.org/10.48550/arXiv.2607.05844}.

\bibitem[\citeproctext]{ref-wu2026agentzero}
Wu, Ming, and Pengyuan Zhu. 2026. {`Agent Zero Memory: Provenance-Aware
Long-Term Memory for {LLM} Agents'}.
\url{https://doi.org/10.48550/arXiv.2608.29606}.

\bibitem[\citeproctext]{ref-xu2026gpm}
Xu, Guodong. 2026. {`Governed Persistent Memory: Source-Bound State
Semantics and Fail-Closed Release for Long-Horizon Agents'}.
\url{https://doi.org/10.48550/arXiv.2608.12476}.

\bibitem[\citeproctext]{ref-xuj2026provenancelaundering}
Xu, Jinghan, Yiyong Xiao, Wanru Shao, Hankai Liu, and Xinjin Li. 2026.
{`Memory Provenance Laundering in {LLM} Agents: A Non-Amplification
Firewall for Persistent Memory'}.
\url{https://doi.org/10.48550/arXiv.2607.29167}.

\bibitem[\citeproctext]{ref-zerhoudi2026compactioncliff}
Zerhoudi, Saber, Jelena Mitrovic, and Michael Granitzer. 2026. {`The
Compaction Cliff in Long-Running {AI} Agent Memory'}.
\url{https://doi.org/10.48550/arXiv.2608.22752}.

\bibitem[\citeproctext]{ref-zhao2026memorepair}
Zhao, Yang, Chengxiao Dai, Mengying Kou, and Yue Xiu. 2026.
{`{MemoRepair}: Barrier-First Cascade Repair in Agentic Memory'}.
\url{https://doi.org/10.48550/arXiv.2605.07242}.

\bibitem[\citeproctext]{ref-zhu2026tiermem}
Zhu, Qiming, Shunian Chen, Rui Yu, Zhehao Wu, and Benyou Wang. 2026.
{`From Lossy to Verified: A Provenance-Aware Tiered Memory for Agents'}.
\url{https://doi.org/10.48550/arXiv.2602.17913}.

\end{CSLReferences}

\clearpage

\section{Appendices}\label{appendices}

The following appendices are part of the expanded working draft and are
separate from its main argument.

\section{Appendix A. Acceptance
Conditions}\label{appendix-a.-acceptance-conditions}

These conditions are acceptance requirements, not statements that the
implementation passes.

\begin{longtable}[]{@{}
  >{\raggedright\arraybackslash}p{(\columnwidth - 2\tabcolsep) * \real{0.5000}}
  >{\raggedright\arraybackslash}p{(\columnwidth - 2\tabcolsep) * \real{0.5000}}@{}}
\toprule\noalign{}
\begin{minipage}[b]{\linewidth}\raggedright
Area
\end{minipage} & \begin{minipage}[b]{\linewidth}\raggedright
Minimum evidence required before confirmation
\end{minipage} \\
\midrule\noalign{}
\endhead
\bottomrule\noalign{}
\endlastfoot
Input separation & A trace shows the deployment process receives no
\texttt{case.claims}, hidden answer values, gold source IDs, fixed
correct event index, or evaluator-derived lineage \\
Claim discovery & An independently labeled set measures omissions and
false candidates from natural requests and exposed drafts \\
Source opening & Randomized event positions and paginated natural logs
are opened through real tools; clipped and mismatched records are
rejected \\
Typed semantics & Tests cover execution/completion, artifact binding,
current-version authority, historical-time scope, revocation, conflict,
and unknown states \\
Archive completeness & Certificates identify their view and replay
coverage; candidate top-k results alone cannot certify absence of
invalidation \\
Freshness & Old unexpired grant followed by revocation is rejected;
same-head missed revocation, scheduled transition, dropped notification,
out-of-order arrival, restart, and head races are exercised \\
Reuse & A no-reopen path records current head/object checks and does not
silently fall back to TTL \\
Host delivery & Payload hash and a supported host receipt demonstrate
actual delivery after the relevant compaction boundary \\
Final decisions & The evaluator separately verifies action behavior,
obligatory claims, source eligibility, abstention, and model
consumption \\
Baseline fidelity & Raw-RAG configuration and delivered text are
auditable; close-method reproductions and approximations are clearly
distinguished \\
Resource accounting & All extraction, verification, acquisition,
metadata, preprocessing, model, and failure costs are retained \\
\end{longtable}

\section{Appendix B. Required Evidence for a Submission-Ready
Paper}\label{appendix-b.-required-evidence-for-a-submission-ready-paper}

The following gaps remain material:

\begin{enumerate}
\def\labelenumi{\arabic{enumi}.}
\tightlist
\item
  \textbf{A deployed-interface implementation.} Eliminate the old
  structured-claim and fixed-location dependencies, implement complete
  lifecycle qualification and head validation, and provide an auditable
  boundary between public runtime data and hidden labels.
\item
  \textbf{Independent freshness validation.} Demonstrate that the known
  counterexample and the additional failure modes in Section 6.6 are
  addressed. Record finite test scope rather than claiming universal
  correctness.
\item
  \textbf{Strong and close comparisons.} Run B under a frozen faithful
  adaptation, assess a governance implementation, and freeze the TierMem
  or sufficiency-only comparison. A self-written approximation is not
  evidence of reproducing another paper.
\item
  \textbf{Natural, independent tasks.} Collect licensed workflows in
  which post-compaction decisions matter, with repository and
  event-family separation. Retire the old 24 fixtures from formal
  estimation.
\item
  \textbf{A justified sample size.} Complete the pilot, power and
  precision analyses, and a substantive justification for the H2 margin.
  Freeze the hierarchy before outcome access.
\item
  \textbf{Paired confirmatory continuations.} Execute the frozen design,
  retain failures, verify actual evidence delivery, and obtain
  independent behavioral and source-based scores.
\item
  \textbf{Incremental mechanism evidence.} Show what qualification
  changes when candidate evidence is identical, and what certified reuse
  changes when the acquisition pipeline is identical.
\item
  \textbf{Total-cost evidence.} Estimate the joint quality and cost
  criterion with all relevant overhead. If the criterion is not met,
  report the limits without replacing the endpoint.
\item
  \textbf{External validity and failure analysis.} Report every planned
  host/model configuration, including adverse or uncertain results, and
  distinguish continuation replay from full natural runs.
\item
  \textbf{An independently reviewable manuscript and artifact.}
  Recompute each table from released or access-controlled records,
  verify every citation and license, resolve authorship and AI
  disclosure, and obtain human approval of the final claims.
\end{enumerate}

Meeting these conditions makes the empirical claim reviewable; it does
not guarantee a positive result or acceptance by a venue. If the
principal increment is absent, a revised paper should state that finding
and its uncertainty rather than retaining an unsupported methods claim.

\section{Appendix C. Provenance of the Existing Project
Observations}\label{appendix-c.-provenance-of-the-existing-project-observations}

The draft's development statements refer to the following inspected
materials. The legacy source revision is
\texttt{49e772bbf60532682771874b7ec96958f4a179e3}.

\begin{longtable}[]{@{}
  >{\raggedright\arraybackslash}p{(\columnwidth - 4\tabcolsep) * \real{0.3333}}
  >{\raggedright\arraybackslash}p{(\columnwidth - 4\tabcolsep) * \real{0.3333}}
  >{\raggedright\arraybackslash}p{(\columnwidth - 4\tabcolsep) * \real{0.3333}}@{}}
\toprule\noalign{}
\begin{minipage}[b]{\linewidth}\raggedright
Material
\end{minipage} & \begin{minipage}[b]{\linewidth}\raggedright
Identifier
\end{minipage} & \begin{minipage}[b]{\linewidth}\raggedright
What it supports
\end{minipage} \\
\midrule\noalign{}
\endhead
\bottomrule\noalign{}
\endlastfoot
Legacy results & Versioned November 26 results report at the pinned
legacy commit & Reported synthetic-fixture counts, timing of
intervention, and limitations \\
Legacy driver & Native evidence-recovery driver at the pinned legacy
commit & Predefined claim list, structured source filter, fixed event
opening, and read-only continuation \\
Legacy gate & Evidence-check gate at the pinned legacy commit &
Local-evidence shortcut, type/status checks, and absence of a required
current archive head \\
Local diagnostic record & \texttt{counterfactual-freshness.json} from
the prior local audit, inspected 28 November 2026 & Recorded
old-grant/later-revocation inputs and the observed no-search output \\
P1 planning document & P1 evidence-verification research plan, an
internal working document & Proposed primary endpoint, statistical
hierarchy, task plan, and resource assumptions \\
Live-state engineering & \texttt{LIVE\_STATE.md} and the corresponding
implementation and regression suite & Bounded current-workspace
observations, execution and claim binding, finite test coverage, and
explicit concurrency limits \\
Experiment controller & \texttt{EXPERIMENT.md} and its controller
contract tests & Local freeze, first-attempt dispatch, evaluator join,
component ledgers, and remaining review boundaries \\
Online development & \texttt{artifacts/live-native/} manifest, attempt
summary, and both extraction records & One approval-blocked integration
attempt, one semantic extraction miss, one same-family prompt retest,
and unmeasured costs \\
Candidate acquisition & \texttt{tasks/acquisition/} manifests and source
records & Sixteen source candidates, eight projects, fixed provenance,
ancestry and label cautions, and unfilled P1 acceptance fields \\
Program reproduction & \texttt{tasks/reproduction/} commands, phase
hashes, stdout, and JUnit reports & Two source/environment checks, four
actual pytest invocations, fixed patches, observed outcomes, and
retained failures \\
\end{longtable}

The local diagnostic is not presented as a publicly archived dataset or
a peer-reviewed source. Before submission it must be included in a
reproducible, appropriately anonymized artifact with its own stable
identifier. Statements about the legacy study are attributed to its
report and inspected code; the entire legacy model experiment was not
rerun for this draft.

\section{Appendix D. Verified Primary Sources and Comparison
Notes}\label{appendix-d.-verified-primary-sources-and-comparison-notes}

Checked on 28 November 2026. This is a focused overlap review, not an
exhaustive systematic review. Full-text sections were inspected for the
principal five requested works and for TierMem, MSS-Complement, and
CompAct-Bench. Additional records below were checked at their primary
abstract/metadata pages; their detailed implementation behavior should
not be inferred beyond that inspected material.

\begin{longtable}[]{@{}
  >{\raggedright\arraybackslash}p{(\columnwidth - 4\tabcolsep) * \real{0.3333}}
  >{\raggedright\arraybackslash}p{(\columnwidth - 4\tabcolsep) * \real{0.3333}}
  >{\raggedright\arraybackslash}p{(\columnwidth - 4\tabcolsep) * \real{0.3333}}@{}}
\toprule\noalign{}
\begin{minipage}[b]{\linewidth}\raggedright
Reference
\end{minipage} & \begin{minipage}[b]{\linewidth}\raggedright
Primary source and inspected version
\end{minipage} & \begin{minipage}[b]{\linewidth}\raggedright
Relevant distinction for this draft
\end{minipage} \\
\midrule\noalign{}
\endhead
\bottomrule\noalign{}
\endlastfoot
Raw-RAG & \href{https://arxiv.org/html/2610.30755v1}{Raw-RAG, v1} &
Strong raw-history retrieval is the primary challenge; full-turn
expansion matters for budgets \\
GPM & \href{https://arxiv.org/abs/2608.12476v2}{GPM, v2 record};
\href{https://arxiv.org/pdf/2608.12476}{full paper} & Source-bound
stable-head release already exists; current claims must distinguish
free-form continuation from deterministic service output \\
HELM & \href{https://aclanthology.org/2026.emnlp-main.721/}{ACL record};
\href{https://aclanthology.org/2026.emnlp-main.721.pdf}{full paper} &
Governance, tool status, versions, provenance, and memory-operation
control directly overlap \\
Compaction-Friendly &
\href{https://arxiv.org/html/2610.31778v2}{Compaction-Friendly, v2} &
Constraint loss despite task performance is established motivation \\
Folio & \href{https://arxiv.org/html/2611.16627v1}{Folio, v1} & Stable
handles, multiresolution retention, and recoverable original events are
established \\
TierMem & \href{https://arxiv.org/html/2602.17913v1}{TierMem, v1} &
Selective escalation from summaries to raw logs is a close prior
method \\
MSS-Complement &
\href{https://arxiv.org/html/2609.20050v3}{MSS-Complement / SERBench,
v3} & State-conditioned sufficient evidence and separation of discovery
from assembly are already explicit \\
CompAct-Bench &
\href{https://arxiv.org/html/2611.25673v1}{CompAct-Bench, v1} &
Controlled continuation and downstream compaction effects are
established evaluation designs \\
Agent Zero Memory & \href{https://arxiv.org/abs/2608.29606v1}{Agent Zero
Memory, v1} & Intent gating and citations restricted to opened evidence
overlap \\
MemIR & \href{https://arxiv.org/abs/2605.25869v1}{MemIR, v1} & Typed
separation of evidence, cues, and claims overlaps \\
MemoRepair & \href{https://arxiv.org/abs/2605.07242v1}{MemoRepair, v1} &
Invalidation and validated republication are established \\
StateFuse & \href{https://arxiv.org/abs/2607.05844v1}{StateFuse, v1} &
Visible conflicts and auditable correction are established \\
Compaction Cliff & \href{https://arxiv.org/abs/2608.22752v1}{Compaction
Cliff, v1} & Type-dependent retention and repeated compaction are prior
topics \\
Provenance laundering &
\href{https://arxiv.org/abs/2607.29167v1}{Provenance laundering, v1} &
Preserving source authority is distinct from verifying historical
authorization facts \\
Eywa & \href{https://arxiv.org/abs/2605.30771v1}{Eywa, v1} & Immutable
source evidence, typed support, and separation of retrieval from answer
policy overlap \\
\end{longtable}

\end{document}
